\section{起诉}
\begin{enumerate}
  \item 分类。
        \begin{enumerate}
          \item 公诉垄断主义，公诉兼自诉制度。我国公诉为主自诉为辅。
          \item 起诉条件。起诉法定主义，起诉便宜主义（裁量）。我国起诉法定主义为主兼采起诉便宜主义（酌定不起诉）。
        \end{enumerate}
  \item 提起公诉的程序。
        \begin{enumerate}
          \item 检察院审查。
                \begin{enumerate}
                  \item 7 天内进行审查；期限：普通 1 月 + 15 天，速裁 10 天 + 5 天。
                  \item 全方面的实质性的审查。事实、证据、犯罪性质、罪名。
                  \item 应当讯问犯罪嫌疑人，听取辩护人、被害人及其诉讼代理人意见。
                  \item 材料不齐备要求\zd{补送}材料；\zd{同案犯}在逃你逃你的我查我的；
                  \item 管辖。
                        \begin{enumerate}
                          \item 门不当户不对移送 + 通知；
                          \item 误查监委，商量；误查公安，查清了诉，没查清就送该谁给谁；
                          \item 公监误查，查清了 + 双方同意。
                        \end{enumerate}
                  \item 应当鉴定/复验复查而没有，原则公安 \hz 监委重新鉴定，必要时检察院可以自行安排。
                  \item \zd{漏罪漏人}。漏罪要查应当要公安补充侦查；漏人要送要公安补充移送；直接提起公诉。
                  \item 发现新罪。新事实该谁给谁；对已查清的应当依法提起公诉。
                  \item 证据核实。
                \end{enumerate}
          \item 起诉。事实清楚，证据充分，应当追究刑事责任 \yjt 起诉书，同级法院提起公诉。
          \item \zd{不起诉种类}。不起诉决定书。
                \begin{enumerate}
                  \item 法定不起诉。\zd{无犯罪事实} + 法定不追责情况。
                        \begin{enumerate}
                          \item 公安/监委过来的案子，同级/上一级检察院检察长批准。
                          \item 对案还不能结束，说明理由，将材料退回到公安/监委重新侦查/调查。
                          \item 自侦部门，退回然后建议撤掉（不打自己人的脸）。
                        \end{enumerate}
                  \item 酌定不起诉。可诉可不诉。\zd{有罪轻微不判刑，可以酌定不起诉}。
                        \begin{enumerate}
                          \item 一定已经构成犯罪；犯罪情节轻微（不需要判刑 \hz 免除处罚）。
                        \end{enumerate}
                  \item 存疑不起诉。针对补侦。
                  \item 附条件不起诉。针对未成年。小孩犯罪 456，1 年以下够起诉，悔过可以附条件。
                \end{enumerate}
          \item 救济。
                \begin{enumerate}
                  \item 公安。复议 \yjt 复核。
                  \item 监察。向上一级检察院复议（因为是上一级批准的，见上）。
                  \item 被害人。上一级检察院申诉；自诉。不分顺序。
                  \item 被不起诉人。仅对\zd{酌定不起诉}决定可以向原机关申诉。
                \end{enumerate}
        \end{enumerate}
  \item 提起自诉的程序。了解即可。同民诉差不多。
\end{enumerate}
